% Provenance (never printed): see sources/outline_sources.md (Year 2: tholeiitic, calc-alkaline
% and alkaline series; the TAS diagram; settings; isotopic signatures of mantle and crust).
% Classification from BGS RR 99-06 (after Streckeisen 1976, Le Bas et al. 1986, Le Maitre et
% al. 1989); series and data from Vermeesch and Pease (2021) and the GeoplotR package;
% curves and counts from figdata/bachelor-2/series.py; exercise and problem values in the
% Book 4 scratch calculators calc/ch09.py and calc/ch09pb.py.
\chapter{المتتاليات الصهارية والأطر التكتونية}\label{ch:b2:magmatic-series}

ناوِلْ عالمةَ صخور تحليلًا كيميائيًا دون خريطة، وكثيرًا ما تستطيع أن تقول هل
ثارت الحمم عند ظهرة، أو على جزيرة محيطية، أو فوق صفيحة مندسّة. فحمم
آيسلندا وحمم براكين الكاسكيد في غرب أمريكا الشمالية تمتدّ كلتاهما
من البازلت إلى الريوليت؛ وإذا حُلّلت بالمئات تبعت مسارين مختلفين وهي
تتطوّر، أحدهما يزداد غنًى بالحديد، والآخر لا. يسمّي هذا الفصل أولًا
الصخور النارية، من معادنها حين تُرى، ومن
كيميائها حين لا تُرى؛ ثم يجمعها في متتاليات، أي عائلات من الصخور
مترابطة بالتمايز، ويربط كل متتالية بإطار تكتوني، مع
المخطّطات التي يستعملها علماء الصخور لقراءة الإطار من تحليل، وحدودها.
% ledger: vp21:series

\begin{recall}
\textbf{من مجلد المرحلة الثانوية}: البازلت والأنديزيت والريوليت والغابرو والغرانيت؛
والصخور الفلسية والمافية؛ والنسج البركانية والجوفية. \textbf{من مجلد السنة
الأولى}: الظهرات وسط المحيطية، ونطاقات الاندساس والأقواس البركانية، والنقاط الساخنة
والأقاليم النارية الكبرى، والصدوع القارية. \textbf{من \cref{ch:b2:elements-in-the-earth}}:
بازلت MORB وبازلت OIB، \omterm{def:b2:elements-in-the-earth:normalised}{والمخطّطات المعايَرة} إلى \omterm{def:b2:elements-in-the-earth:bse}{الوشاح البدائي} (المخطّطات العنكبوتية)، والعناصر المتوافقة
وغير المتوافقة. \textbf{من \cref{ch:b2:isotope-geochemistry}}: \omterm{def:b2:isotope-geochemistry:reservoirs}{الوشاح
المستنزَف} والخزّانات النظائرية المُثراة. \textbf{من
\cref{ch:b2:magmatic-differentiation}}: \omterm{def:b2:magmatic-differentiation:crystallisation}{التبلور التجزيئي}، \omterm{def:b2:magmatic-differentiation:lld}{وخط النزول
السائل}، \omterm{def:b2:magmatic-differentiation:harker}{ومخطّطات هاركر}، وقانون رايلي.
\end{recall}

\section{تسمية الصخور النارية}\label{sec:b2:magmatic-series:naming}

الصخر الخشن الحبيبات يُظهر معادنه، ويُسمّى بها. أمّا الحمم فكثيرًا ما تكون أدقّ
حبيبات أو أكثر زجاجية من ذلك، فتُسمّى من تحليلها الكيميائي. والنظامان كلاهما
اصطلاحان دوليان: يصفان، ولا يقتضيان شيئًا عن كيفية
تكوّن الصخر.
% ledger: bgs:tas-use

\begin{definition}[التركيب المعدني النمطي؛ مخطّط QAPF]\label{def:b2:magmatic-series:qapf}
\emph{التركيب المعدني النمطي}\index{التركيب المعدني النمطي} لصخر هو نسبة كل معدن من معادنه
حجمًا. و\emph{مخطّط QAPF}\index{مخطّط QAPF} يسمّي
صخرًا ناريًا خشن الحبيبات فيه أقل من \qty{90}{\percent} من المعادن المافية، انطلاقًا من
أربعة مكوّنات نمطية يُعاد حسابها بحيث يكون مجموعها 100: الكوارتز (Q)، وفلسبار القلويات
(A)، والبلاجيوكليز (P)، والفلسباتويدات (F). وبما أن الكوارتز والفلسباتويدات لا
يتبلوران معًا، يقع كل صخر في أحد مثلثين، QAP أو FAP، يلتقيان
على امتداد ضلعهما المشترك AP.
\end{definition}
% ledger: bgs:qapf

\begin{omfigure}
\centering
\begin{tikzpicture}[font=\scriptsize, x=1cm, y=1cm]
\def\L{6.4}
\pgfmathsetmacro{\H}{\L*0.8660254}
% upper point (q, p): q = Q fraction, p = P/(A+P)
\newcommand{\qapfup}[2]{({(#1)*\L/2 + (#2)*(1-(#1))*\L}, {(#1)*\H})}
\newcommand{\dn}[2]{({(#1)*\L/2 + (#2)*(1-(#1))*\L}, {-(#1)*\H})}
\draw[thick] (0,0) -- (\L/2,\H) -- (\L,0) -- (\L/2,-\H) -- cycle;
\draw (0,0) -- (\L,0);
\foreach \q in {0.9,0.6,0.2,0.05} \draw \qapfup{\q}{0} -- \qapfup{\q}{1};
\foreach \q in {0.1,0.6} \draw \dn{\q}{0} -- \dn{\q}{1};
\foreach \p in {0.1,0.65,0.9} \draw \qapfup{0.6}{\p} -- \qapfup{0}{\p};
\draw[dashed] \qapfup{0.6}{0.35} -- \qapfup{0.2}{0.35};
\draw \qapfup{0.2}{0.35} -- \qapfup{0}{0.35};
\foreach \p in {0.1,0.35,0.65,0.9} \draw \dn{0}{\p} -- \dn{0.1}{\p};
\foreach \p in {0.1,0.5,0.9} \draw \dn{0.1}{\p} -- \dn{0.6}{\p};
\node[anchor=south] at (\L/2,\H) {\normalsize Q};
\node[anchor=east] at (0,0) {\normalsize A};
\node[anchor=west] at (\L,0) {\normalsize P};
\node[anchor=north] at (\L/2,-\H) {\normalsize F};
\node at \qapfup{0.94}{0.5} {1a};
\node at \qapfup{0.75}{0.5} {1b};
\node[font=\tiny] at \qapfup{0.4}{0.05} {2};
\node at \qapfup{0.4}{0.225} {3a};
\node at \qapfup{0.4}{0.5} {3b};
\node at \qapfup{0.4}{0.775} {4};
\node[font=\tiny] at \qapfup{0.4}{0.95} {5};
\foreach \p/\n in {0.225/7*,0.5/8*,0.775/9*} \node at \qapfup{0.125}{\p} {\n};
\node[font=\tiny] at \qapfup{0.125}{0.05} {6*};
\node[font=\tiny] at \qapfup{0.125}{0.95} {10*};
\foreach \p/\n in {0.225/7,0.5/8,0.775/9} \node at \qapfup{0.025}{\p} {\n};
\foreach \p/\n in {0.225/7',0.5/8',0.775/9'} \node at \dn{0.05}{\p} {\n};
\node[font=\tiny] at \dn{0.3}{0.05} {11};
\node at \dn{0.35}{0.3} {12};
\node at \dn{0.35}{0.7} {13};
\node[font=\tiny] at \dn{0.3}{0.95} {14};
\node at \dn{0.75}{0.5} {15};
\node[anchor=east] at \qapfup{0.9}{0} {90};
\node[anchor=east] at \qapfup{0.6}{0} {60};
\node[anchor=east] at \qapfup{0.2}{0} {20};
\node[anchor=east] at \dn{0.1}{0} {10};
\node[anchor=east] at \dn{0.6}{0} {60};
\foreach \p/\n in {0.1/10,0.35/35,0.65/65,0.9/90} \node[anchor=south, fill=white, inner sep=0.5pt] at \qapfup{0.6}{\p} {\tiny\n};
\node[anchor=west, align=left, font=\scriptsize] at (\L+0.9,1.6) {\foreignlanguage{arabic}{المثلث العلوي: نسبة Q}\\\foreignlanguage{arabic}{في Q + A + P؛}\\
  \foreignlanguage{arabic}{خطوط القيمة الثابتة}\\\foreignlanguage{arabic}{للمقدار $100\,P/(A+P)$}\\\foreignlanguage{arabic}{عند 10 و35 و65 و90}};
\node[anchor=west, align=left, font=\scriptsize] at (\L+0.9,-1.7) {\foreignlanguage{arabic}{المثلث السفلي: نسبة F}\\\foreignlanguage{arabic}{في F + A + P؛}\\
  \foreignlanguage{arabic}{المقدار $100\,P/(A+P)$}\\\foreignlanguage{arabic}{عند 10 و50 و90}};
\end{tikzpicture}
\par\smallskip{\small تصنيف QAPF للصخور النارية الخشنة الحبيبات. 1a
الكوارتزوليت، 1b الغرانيتويد الغني بالكوارتز، 2 غرانيت فلسبار القلويات، 3a السيانوغرانيت
و3b المونزوغرانيت (وهما معًا الغرانيت)، 4 الغرانوديوريت، 5 التوناليت؛ ومن 6* إلى 10*
الأصناف الحاوية للكوارتز من الحقول 6 إلى 10: 6 سيانيت فلسبار القلويات، 7 السيانيت، 8 المونزونيت، 9 المونزوديوريت
أو المونزوغابرو، 10 الديوريت أو الغابرو أو الأنورثوزيت؛ ومن 6' إلى 10' أصنافها الحاوية
للفلسباتويدات؛ 11 سيانيت الفلسباتويدات، 12 مونزوسيانيت الفلسباتويدات، 13 مونزوديوريت الفلسباتويدات أو
مونزوغابرو الفلسباتويدات، 14 ديوريت الفلسباتويدات أو غابرو الفلسباتويدات، 15 الفويدوليت. والحقول الرقيقة 6 و6' و10 و10'
على الحواف غير مرقّمة. عن نظام هيئة المسح الجيولوجي البريطانية (المبني
على شترايكايزن 1976).}
\end{omfigure}
% ledger: bgs:qapf

يضمّ الحقل 10 ثلاثة أسماء. \omterm{def:b2:magmatic-series:qapf}{فمخطّط QAPF} لا يفصل بين الديوريت والغابرو
والأنورثوزيت، ولا بين الصخور الغابروية فيما بينها؛ بل تفعل ذلك مخطّطات أخرى، مبنيّة على
البلاجيوكليز والمعادن المافية.

\begin{definition}[مخطّط مجموع القلويات--السيليكا]\label{def:b2:magmatic-series:tas}
\emph{مخطّط مجموع القلويات--السيليكا}\index{مخطّط مجموع القلويات--السيليكا} (مخطّط
TAS) يسمّي صخرًا بركانيًا لا يمكن قياس تركيبه المعدني النمطي، من نسبتيه المئويتين
الكتليتين من \ce{Na2O + K2O} بدلالة \ce{SiO2}، بعد إعادة حساب التحليل
إلى \qty{100}{\percent} خاليًا من الماء وثاني أكسيد الكربون. وهو مقسوم
إلى خمسة عشر حقلًا بخطوط مستقيمة لتقاطعاتها إحداثيات ثابتة
(الشكل بعد الفقرة التالية).
\end{definition}
% ledger: bgs:tas, bgs:tas-use

للمخطّط ثلاث مجموعات من الحقول. ففي الأسفل، عند القلويات المنخفضة، تقع
أسماء المتتالية الشائعة، من البيكروبازلت والبازلت إلى الريوليت، مفصولة عند 45 و52
و57 و\qty{63}{\percent} من السيليكا، وبخط مائل بالنسبة إلى الريوليت. وفوقها تضمّ حقول S والتراكيت صخورًا أغنى
بالقلويات عند السيليكا نفسها؛ وتضمّ حقول U والفونوليت والفويديت صخورًا بلغت من الغنى
بالقلويات، ومن الفقر بالسيليكا، حدًّا تبلور معه الفلسباتويدات. وبعض الحقول
مقسوم بمعيار ثانٍ: يُسمّى التراكيبازلت هاوايت إذا كان $\ce{Na2O} - 2
\ge \ce{K2O}$، وتراكيبازلت بوتاسيًا في غير ذلك؛ والتيفريت الذي فيه أكثر من
\qty{10}{\percent} من الأوليفين في معياره (انظر أدناه) هو بازانيت؛ والريوليت أو التراكيت الذي
تتجاوز فيه النسبة الجزيئية $(\ce{Na2O + K2O})/\ce{Al2O3}$ الواحد يُسمّى فوق قلوي. ويشترط النظام
تحاليل فيها أقل من \qty{2}{\percent} من الماء و\qty{0.5}{\percent} من
ثاني أكسيد الكربون: فالتبدّل يضيف الماء ويحرّك القلويات، \omterm{def:b2:magmatic-differentiation:cumulate}{والمتراكم} ليس
سائلًا.
% ledger: bgs:tas, bgs:tas-use

\begin{omfigure}
\centering
\begin{tikzpicture}
\begin{axis}[omrecord, width=0.86\linewidth, height=8.6cm,
  xmin=37, xmax=80, ymin=0, ymax=16, xtick={37,41,...,77}, ytick={0,2,...,16},
  xlabel={\mbox{\textdir TLT \ce{SiO2} (\unit{\percent})}}, ylabel={\mbox{\textdir TLT \ce{Na2O + K2O} (\unit{\percent})}},
  legend style={font=\scriptsize, at={(0.98,0.98)}, anchor=north east},
  tick label style={font=\scriptsize}, label style={font=\small}]
\addplot[only marks, mark=*, mark size=0.55pt, blue!70!black, opacity=0.6] table[x=SiO2, y=alk] {figdata/out/bachelor-2-series-cath-th.dat};
\addlegendentry{آيسلندا}
\addplot[only marks, mark=*, mark size=0.55pt, red!75!black, opacity=0.6] table[x=SiO2, y=alk] {figdata/out/bachelor-2-series-cath-ca.dat};
\addlegendentry{الكاسكيد}
\foreach \f in {Pc,B,O1,O2,O3,R,S1,S2,S3,T,U1,U2,U3,Ph,F}
  \addplot[black, thin, forget plot] table[x=SiO2, y=alk] {figdata/out/bachelor-2-series-tas-\f.dat};
\node[font=\tiny] at (axis cs:43,1.6) {Pc};
\node[font=\tiny, fill=white, inner sep=0.5pt] at (axis cs:50.5,0.8) {B};
\node[font=\tiny] at (axis cs:54.5,2.2) {O1};
\node[font=\tiny] at (axis cs:60,2.2) {O2};
\node[font=\tiny] at (axis cs:66.5,2.2) {O3};
\node[font=\tiny] at (axis cs:74.5,11) {R};
\node[font=\tiny, fill=white, inner sep=0.5pt] at (axis cs:49,5.8) {S1};
\node[font=\tiny] at (axis cs:52.5,7.4) {S2};
\node[font=\tiny] at (axis cs:57.5,8.9) {S3};
\node[font=\tiny] at (axis cs:64.5,12) {T};
\node[font=\tiny] at (axis cs:44,6.6) {U1};
\node[font=\tiny] at (axis cs:48.7,9.4) {U2};
\node[font=\tiny] at (axis cs:52.5,11.8) {U3};
\node[font=\tiny] at (axis cs:57,14.6) {Ph};
\node[font=\tiny] at (axis cs:39,10) {F};
\end{axis}
\end{tikzpicture}
\par\smallskip{\small مخطّط TAS (عن لو با وآخرين، 1986، كما تقدّمه هيئة
المسح الجيولوجي البريطانية) مع 456 من الحمم والصخور المتداخلة من آيسلندا (بالأزرق)
و580 من الكاسكيد (بالأحمر)، والأكاسيد كما حُلّلت. Pc البيكروبازلت، B البازلت، O1 الأنديزيت
البازلتي، O2 الأنديزيت، O3 الداسيت، R الريوليت؛ S1 التراكيبازلت، S2 التراكيأنديزيت
البازلتي، S3 التراكيأنديزيت، T التراكيت أو التراكيداسيت؛ U1 التيفريت أو
البازانيت، U2 الفونوتيفريت، U3 التيفريفونوليت، Ph الفونوليت، F الفويديت.}
\end{omfigure}
% ledger: bgs:tas, geoplotr:tas, data:b4-cath


\begin{definition}[التشبّع بالسيليكا]\label{def:b2:magmatic-series:saturation}
تكون \omterm{def:b2:magma-emplacement:magma}{الصهارة} أو الصخر \emph{فوق مشبّعة بالسيليكا}\index{التشبّع بالسيليكا} إذا كان فيها
من السيليكا أكثر مما تستطيع أكاسيدها الأخرى أن تتّحد معه في المعادن السيليكاتية، فيظهر
الكوارتز (أو معدن سيليكا آخر)؛ و\emph{دون مشبّعة بالسيليكا} إذا كان فيها
أقل مما يلزم لتحويل كل قلوياتها إلى فلسبار، فيظهر بدلًا منه فلسباتويد (النيفلين،
اللوسيت)، ولا يظهر مع الكوارتز أبدًا؛ و\emph{مشبّعة بالسيليكا}
بين الحالتين.
\end{definition}

وللحمم التي لا يمكن عدّ معادنها، يُحكم على التشبّع من خلال \emph{معيار}: أي
حساب قياسي يوزّع أكاسيد التحليل على معادن مثالية
بترتيب ثابت، القلويات في الفلسبار أولًا، ويعطي فائض السيليكا
كوارتزًا أو نقصها فلسباتويدًا. والمعيار المعتاد، معيار CIPW، يفصل بين
نوعي البازلت: فالبازلت الذي في معياره نيفلين بازلت قلوي؛ والذي
ليس فيه نيفلين بازلت دون قلوي، ومن بين هذه تحمل البازلتات الثولييتية
الهيبرستين (الأورثوبيروكسين) في المعيار، وكثيرًا ما تحمل البيروكسين الفقير بالكالسيوم في التركيب النمطي.
ويطبّق \cref{exo:b2:magmatic-series:11} نسخة مبسّطة منه.
% ledger: bgs:tas-use, bgs:series-terms

\section{المتتاليات الصهارية}\label{sec:b2:magmatic-series:series}

قلّما تتشابه حمم بركان واحد، أو إقليم واحد، لكنها مترابطة: فهي
تشتقّ من صهارات أمّ متشابهة بتمايز
\cref{ch:b2:magmatic-differentiation}، وترسم معًا خط نزول.

\begin{definition}[المتتالية الصهارية]\label{def:b2:magmatic-series:series}
\emph{المتتالية الصهارية}\index{المتتالية الصهارية} مجموعة من الصخور النارية من دورة
صهارية واحدة أو إقليم واحد، مترابطة بالتمايز من صهارات أمّ متشابهة
وتشترك في سمات كيميائية ومعدنية مميّزة. وتُميَّز ثلاث منها عادةً. \emph{المتتالية الثولييتية}\index{المتتالية الثولييتية} تبدأ من
بازلتات مشبّعة بالسيليكا أو فوق مشبّعة بها قليلًا، في معيارها الهيبرستين،
وتُظهر إثراءً واضحًا بالحديد في بداية تمايزها. و\emph{المتتالية
الكلسية القلوية}\index{المتتالية الكلسية القلوية} تمتدّ من البازلت عبر
الأنديزيت إلى الداسيت والريوليت، مع إثراء ضئيل بالحديد أو بلا إثراء، ومع
معادن مائية (الهورنبلند، البيوتيت) في أفرادها المتطوّرة. و\emph{المتتالية
القلوية}\index{المتتالية القلوية} تبدأ من بازلتات في معيارها نيفلين وتتطوّر
عبر حقلي S وU في مخطّط TAS نحو التراكيت والفونوليت.
\end{definition}
% ledger: bgs:series-terms, bgs:tas-use

تصف الأسماء مجموعات، لا صخورًا مفردة: فالبازلت المعزول يمكن أن ينتمي إلى أيّ
منها، والأسماء نفسها استُعملت بمعانٍ مختلفة، فعلى عالمة الصخور
أن تقول ماذا تعني بها. والتمييز بين المتتاليتين الثولييتية
والكلسية القلوية يقوم على الحديد، والمخطّط الذي يُظهره أفضل إظهار يرسم ثلاث
مجموعات من الأكاسيد.
% ledger: bgs:series-terms

\begin{definition}[مخطّط AFM]\label{def:b2:magmatic-series:afm}
\emph{مخطّط AFM}\index{مخطّط AFM} مخطّط مثلثي رؤوسه
$A = \ce{Na2O + K2O}$ و$F = \ce{FeO_T}$ (الحديد كله على هيئة \ce{FeO}) و$M = \ce{MgO}$، معبّرًا عنها
بالنسب المئوية الكتلية بعد إعادة حسابها بحيث $A + F + M = 100$. وخطٌّ مرسوم عليه يفصل
الحقل الثولييتي، من جهة F، عن الحقل الكلسي القلوي.
\end{definition}
% ledger: vp21:series, geoplotr:afm

البازلت الذي يتمايز يفقد المغنيسيوم أسرع من أيّ شيء آخر، لأن الأوليفين
والبيروكسين أغنى بالمغنيسيوم من السائل، ويحتفظ بقلوياته، وهي
غير متوافقة: فيبتعد خط نزوله عن M، نحو A. وما يحدث
للحديد يحسم المتتالية. ففي \omterm{def:b2:magma-emplacement:magma}{الصهارة} الثولييتية لا يغادر الحديد إلا مع
السيليكات الحديدوزية الغنية بالمغنيسيوم، وأبطأ من المغنيسيوم، فترتفع النسبة
$F/M$ ويصعد الاتجاه أولًا نحو F قبل أن ينعطف إلى A: إنه
\emph{اتجاه فينر}. وفي \omterm{def:b2:magma-emplacement:magma}{الصهارة} الكلسية القلوية، الأكثر تأكسدًا، يغادر الحديد أيضًا في
الماغنتيت، الذي يحمل الحديد الحديديك، فتبقى $F/M$ شبه ثابتة بينما
يتّجه السائل مباشرة نحو A: إنه \emph{اتجاه بوين}. ويُظهر الشكل أدناه الاتجاهين في
صخور حقيقية، ثولييتات آيسلندا وحمم الكاسكيد الكلسية القلوية. أمّا
الخط الفاصل، المرسوم بالعين عام 1971، فيفصل بينهما جيدًا حيث تكون بازلتية؛ وعند
القلويات المرتفعة يتقارب الاتجاهان ولا يقول الخط شيئًا يُذكر.
% ledger: vp21:series, geoplotr:afm, data:b4-cath

\begin{omfigure}
\centering
\begin{tikzpicture}
\begin{axis}[name=afm, width=0.5\linewidth, height=0.5\linewidth, axis equal image,
  hide axis, xmin=-0.04, xmax=1.04, ymin=-0.06, ymax=0.92, clip=false]
\addplot[only marks, mark=*, mark size=0.5pt, blue!70!black, opacity=0.6] table[x=afm_x, y=afm_y] {figdata/out/bachelor-2-series-cath-th.dat};
\addplot[only marks, mark=*, mark size=0.5pt, red!75!black, opacity=0.6] table[x=afm_x, y=afm_y] {figdata/out/bachelor-2-series-cath-ca.dat};
\addplot[very thick, black] table[x=x, y=y] {figdata/out/bachelor-2-series-afm-line.dat};
\draw[thick] (axis cs:0,0) -- (axis cs:1,0) -- (axis cs:0.5,0.8660254) -- cycle;
\node[anchor=south] at (axis cs:0.5,0.866) {F};
\node[anchor=north east] at (axis cs:0,0) {A};
\node[anchor=north west] at (axis cs:1,0) {M};
\node[font=\scriptsize, fill=white, inner sep=1pt] at (axis cs:0.52,0.62) {ثولييتي};
\node[font=\scriptsize, fill=white, inner sep=1pt] at (axis cs:0.33,0.12) {كلسي قلوي};
\end{axis}
\begin{axis}[omrecord, at={(afm.east)}, anchor=west, xshift=1.3cm,
  width=0.46\linewidth, height=0.46\linewidth,
  xmin=-2.5, xmax=2.2, ymin=-4.5, ymax=1.5,
  xlabel={$\ln(A/F)$}, ylabel={$\ln(M/F)$},
  legend style={font=\scriptsize, at={(0.98,0.98)}, anchor=north east},
  tick label style={font=\scriptsize}, label style={font=\small}]
\addplot[only marks, mark=*, mark size=0.5pt, blue!70!black, opacity=0.6] table[x=lnAF, y=lnMF] {figdata/out/bachelor-2-series-lr-th.dat};
\addlegendentry{آيسلندا}
\addplot[only marks, mark=*, mark size=0.5pt, red!75!black, opacity=0.6] table[x=lnAF, y=lnMF] {figdata/out/bachelor-2-series-lr-ca.dat};
\addlegendentry{الكاسكيد}
\addplot[thick, blue!50!black, forget plot] table[x=x, y=y_th] {figdata/out/bachelor-2-series-lr-fits.dat};
\addplot[thick, red!50!black, forget plot] table[x=x, y=y_ca] {figdata/out/bachelor-2-series-lr-fits.dat};
\end{axis}
\end{tikzpicture}
\par\smallskip{\small اليسار: \omterm{def:b2:magmatic-series:afm}{مخطّط AFM} لصخور آيسلندا (بالأزرق) والكاسكيد (بالأحمر)
في شكل TAS، مع الخط الفاصل لإرفاين وباراغار (1971). اليمين:
الصخور نفسها في إحداثيات لوغاريتم النسب، حيث تصبح كل متتالية شريطًا شبه
مستقيم؛ والخطان انحداران متعامدان، ميلاهما $-1.10$ (آيسلندا) و$-0.72$
(الكاسكيد)، يلتقيان عند $(-3.84, 2.12)$ كما وجد فيرميش وبيز (2021).}
\end{omfigure}
% ledger: geoplotr:afm, data:b4-cath, vp21:model

من صخور آيسلندا الواقعة ضمن مدى الخط الفاصل، يقع 351 في
الجهة الثولييتية و23 في الأخرى؛ ومن صخور الكاسكيد، 417 في الجهة الكلسية القلوية
و57 في الأخرى. والمنحنى حدّ أخرق، ونموذج بسيط يفسّر
شكل الاتجاهات. لنفترض أن \omterm{def:b2:magma-emplacement:magma}{الصهارة} تفقد قلوياتها وحديدها ومغنيسيومها
بمعدلات متناسبة مع الكميات التي تحويها، $\mathrm dX/\mathrm dt = -\lambda_X X$.

\begin{proposition}[اتجاهات لوغاريتم النسب]\label{prop:b2:magmatic-series:logratio}
إذا تناقصت $A$ و$F$ و$M$ أُسّيًا بالثوابت $\lambda_A$ و$\lambda_F$
و$\lambda_M$، فإن \omterm{def:b2:magma-emplacement:magma}{الصهارة} تتحرّك على خط مستقيم في مستوى $\ln(A/F)$
و$\ln(M/F)$، ميله $(\lambda_F - \lambda_M)/(\lambda_F - \lambda_A)$، بينما تتحرّك
على منحنٍ في مثلث AFM.
\end{proposition}

\begin{proof}
مع $X = X_0 e^{-\lambda_X t}$، $\ln(A/F) = \ln(A_0/F_0) + (\lambda_F - \lambda_A)t$
و$\ln(M/F) = \ln(M_0/F_0) + (\lambda_F - \lambda_M)t$: كلاهما خطي في $t$، ومن ثمّ يكون
أحدهما خطيًا في الآخر، بالميل المذكور. والنسبتان $A/F$ و$M/F$ تحدّدان
النقطة على المثلث، التي إحداثياتها $A/(A + F + M)$، \dots\ ليست خطية في $t$.
\end{proof}

الميل الأشدّ لصخور آيسلندا يعني $\lambda_F$ أصغر: فالحديد يغادر
\omterm{def:b2:magma-emplacement:magma}{الصهارة} الثولييتية أبطأ (\cref{exo:b2:magmatic-series:10}). وقد بيّن فيرميش وبيز أن
شريطي الشكل أعلاه يلتقيان في نقطة واحدة، قرآها مصدرًا مشتركًا
افتراضيًا.
% ledger: vp21:model

تختلف المتتاليتان الثولييتية والكلسية القلوية في بوتاسها أيضًا. فمخطّط
\ce{K2O}--\ce{SiO2} يقسم البازلتات والأنديزيتات والداسيتات والريوليتات التي فيها
أكثر من \qty{48}{\percent} من السيليكا إلى أنماط منخفضة البوتاسيوم ومتوسطة البوتاسيوم وعالية البوتاسيوم بخطّين
مستقيمين (الشكل أدناه). ومعظم المجموعتين هنا متوسط البوتاسيوم: فالمخطّط
يقول كم من البوتاس تحمل متتالية، لا أيّ متتالية هي.
% ledger: bgs:k2o-sio2, data:b4-cath

\begin{omfigure}
\centering
\begin{tikzpicture}
\begin{axis}[omrecord, width=0.78\linewidth, height=6.2cm,
  xmin=45, xmax=79, ymin=0, ymax=5,
  xlabel={\mbox{\textdir TLT \ce{SiO2} (\unit{\percent})}}, ylabel={\mbox{\textdir TLT \ce{K2O} (\unit{\percent})}},
  legend style={font=\scriptsize, at={(0.02,0.98)}, anchor=north west},
  tick label style={font=\scriptsize}, label style={font=\small}]
\addplot[only marks, mark=*, mark size=0.55pt, blue!70!black, opacity=0.6] table[x=SiO2, y=K2O] {figdata/out/bachelor-2-series-cath-th.dat};
\addlegendentry{آيسلندا}
\addplot[only marks, mark=*, mark size=0.55pt, red!75!black, opacity=0.6] table[x=SiO2, y=K2O] {figdata/out/bachelor-2-series-cath-ca.dat};
\addlegendentry{الكاسكيد}
\addplot[thick, black, forget plot] table[x=x, y=y_low] {figdata/out/bachelor-2-series-k2o-lines.dat};
\addplot[thick, black, forget plot] table[x=x, y=y_high] {figdata/out/bachelor-2-series-k2o-lines.dat};
\draw[thick] (axis cs:48,0) -- (axis cs:48,0.3);
\node[font=\scriptsize, fill=white, inner sep=1pt] at (axis cs:60,3.7) {عالي البوتاسيوم};
\node[font=\scriptsize, fill=white, inner sep=1pt] at (axis cs:74.8,2.2) {متوسط البوتاسيوم};
\node[font=\scriptsize, fill=white, inner sep=1pt] at (axis cs:60,0.35) {منخفض البوتاسيوم};
\end{axis}
\end{tikzpicture}
\par\smallskip{\small البوتاس بدلالة السيليكا لصخور آيسلندا والكاسكيد، مع
الخطّين المارّين بالنقاط \babelsublr{(48, 0.3)--(68, 1.2)} و\babelsublr{(48, 1.2)--(68, 3.1)} اللذين يفصلان الأنماط المنخفضة البوتاسيوم
والمتوسطة البوتاسيوم والعالية البوتاسيوم (عن بيتشيريلو وتايلور، 1976، كما اعتمدهما
النظام الدولي)، ممدودين هنا إلى ما بعد 68.}
\end{omfigure}
% ledger: bgs:k2o-sio2, data:b4-cath

لماذا يختلف التأكسد؟ تتكوّن الصهارات الكلسية القلوية فوق نطاقات الاندساس، حيث
يمدّ اللوحُ الذي يفقد ماءه إسفينَ الوشاح بأنواع كيميائية مؤكسدة؛ فتبلور الصهارات
الماغنتيت، وأفرادها المتطوّرة الهورنبلند أو البيوتيت. أمّا الصهارات الثولييتية
فتتكوّن عند الظهرات، حيث ينصهر الوشاح بتخفيف الضغط
(\cref{ch:b2:mantle-melting})، وهي أكثر اختزالًا.
% ledger: vp21:series

\section{المتتاليات والأطر}\label{sec:b2:magmatic-series:settings}

كل إطار في تكتونية الصفائح ينتج صهاراته الخاصة، لأن كلًّا منها يصهر
الوشاح بطريقته: بتخفيف الضغط تحت الظهرات والنقاط الساخنة، وبإضافة الماء
فوق نطاقات الاندساس (\cref{ch:b2:mantle-melting}).

\emph{الظهرات وسط المحيطية} تقذف بازلتات ثولييتية، هي بازلت MORB في
\cref{ch:b2:elements-in-the-earth}، فقيرة بالعناصر غير المتوافقة وتحمل
البصمة النظائرية المستنزَفة لمصدرها. وتضيف آيسلندا، وهي ظهرة فوق نقطة ساخنة،
ثولييتاتها إلى تلك القائمة، مع متتالية كاملة حتى الريوليت.
% ledger: vp21:series

\emph{الجزر المحيطية} تنمو على مراحل. فبركان هاوايي يقذف حممًا قلوية وهو
فتيّ، ثم يبني درعه من البازلت الثولييتي، أي ما يقدَّر بنسبة من 80 إلى
\qty{95}{\percent} من حجمه النهائي، في أول مليون إلى مليوني سنة من عمره؛ وحين
ينجرف بعيدًا عن النقطة الساخنة تعود حممه قلوية، منتهيةً بذيل طويل من
الثورانات القلوية النادرة. فالانصهار الأقل يعطي صهارات أكثر قلوية: لأن القلويات،
لكونها غير متوافقة، تُخفَّف أقل في كسر مصهور صغير. ونظائريًا، تأخذ بازلتات
الجزر المحيطية عيّنات من الخزّانات المُثراة في \cref{ch:b2:isotope-geochemistry}.
% ledger: usgs14:hawaii-stages

\emph{الأقواس البركانية} تقذف متتاليات كلسية قلوية. وتحمل عناصرها النزرة
علامة مميّزة (الشكل أدناه): فبعد المعايرة إلى \omterm{def:b2:elements-in-the-earth:bse}{الوشاح البدائي}، تكون بازلتات
الأقواس مُثراة بعناصر Rb وBa وTh وU وPb وSr، وهي عناصر يضيفها اللوح إلى
الإسفين، ومستنزَفة بالنيوبيوم والتنتالوم مقارنةً بجيرانهما. ولبازلت
الأقواس الوسيط في مجموعة البيانات $\mathrm{Nb}/\mathrm{La} = 0.49$، مقابل 0.97 لبازلت MORB
و1.23 لبازلتات الجزر المحيطية.
% ledger: data:b4-basalt-settings

\begin{omfigure}
\centering
\begin{tikzpicture}
\begin{axis}[omrecord, width=0.86\linewidth, height=6.4cm, ymode=log,
  xmin=0.5, xmax=22.5, ymin=1, ymax=300,
  xtick={1,...,22}, xticklabels={Rb,Ba,Th,U,Nb,Ta,La,Ce,Pb,Sr,Nd,Zr,Hf,Sm,Eu,Ti,Gd,Dy,Y,Er,Yb,Lu},
  ylabel={\foreignlanguage{arabic}{العيّنة / الوشاح البدائي}},
  legend style={font=\scriptsize, at={(0.98,0.98)}, anchor=north east},
  tick label style={font=\scriptsize}, label style={font=\small}]
\addplot[thick, mark=square*, mark size=1.4pt, orange!85!black] table[x=idx, y=oib] {figdata/out/bachelor-2-series-spider.dat};
\addlegendentry{بازلت الجزر المحيطية}
\addplot[thick, mark=*, mark size=1.4pt, blue!70!black] table[x=idx, y=morb] {figdata/out/bachelor-2-series-spider.dat};
\addlegendentry{MORB}
\addplot[thick, mark=triangle*, mark size=1.7pt, red!75!black] table[x=idx, y=iab] {figdata/out/bachelor-2-series-spider.dat};
\addlegendentry{بازلت الأقواس الجزرية}
\end{axis}
\end{tikzpicture}
\par\smallskip{\small التراكيب الوسيطة المعايَرة إلى \omterm{def:b2:elements-in-the-earth:bse}{الوشاح البدائي} لعدد 241 بازلت MORB، و259
بازلت جزر محيطية و256 بازلت أقواس جزرية (مجموعة فيرميش، 2006؛
\omterm{def:b2:elements-in-the-earth:bse}{والوشاح البدائي} لماكدونو وصن، 1995). تُظهر بازلتات الأقواس الحوض عند Nb
وTa والقمم عند Ba وPb وSr؛ وبازلتات الجزر المحيطية مُثراة بكل
العناصر الأشدّ عدم توافق.}
\end{omfigure}
% ledger: data:b4-basalt-settings, ms95:pm-trace

داخل القارات تلاقي \omterm{def:b2:magma-emplacement:magma}{الصهارة} البازلتية قشرة سميكة تستطيع أن تسخّنها وتصهرها،
وقد تكون النتيجة نشاطًا بركانيًا ذا نمطين.

\begin{definition}[النشاط البركاني ثنائي النمط]\label{def:b2:magmatic-series:bimodal}
\emph{النشاط البركاني ثنائي النمط}\index{النشاط البركاني ثنائي النمط} هو ثوران مجموعتين من الصهارات، في إقليم واحد،
بازلتية وريوليتية، مع قليل من التراكيب الوسيطة
بينهما.
\end{definition}

النشاط البركاني على امتداد مسار النقطة الساخنة يلوستون ثنائي النمط: ففي سهل
نهر سنيك الشرقي، تغطّي بازلتاتٌ أحدث فيها نحو \qty{50}{\percent} من السيليكا ريوليتاتٍ أقدم فيها نحو \qty{75}{\percent}.
والفجوة بينهما هي المفتاح: فتمايز البازلت وحده كان سيخلّف حممًا
وسيطة، والفجوة تشير إلى مصدر ثانٍ للريوليتات، كالقشرة التي يسخّنها البازلت.
% ledger: yvo23:bimodal

\emph{الأقاليم النارية الكبرى}، التي مرّت في مجلد السنة الأولى، تقذف البازلت بأحجام
لا يضاهيها إطار آخر؛ وتُقرأ متتاليتها بالمخطّطات نفسها.

\section{مخطّطات التمييز}\label{sec:b2:magmatic-series:discrimination}

توحي أنماط القسم السابق بقراءة الإطار من التحليل
مباشرة، وقد رسم علماء الصخور مخطّطات كثيرة لذلك.

\begin{definition}[مخطّط التمييز]\label{def:b2:magmatic-series:discrimination}
\emph{مخطّط التمييز}\index{مخطّط التمييز} مخطّط لمتغيّرات
كيميائية، عادةً عناصر نزرة غير متحرّكة، مقسوم إلى حقول مرسومة
حول تراكيب صخور من أطر معروفة، يُستعمل لنسبة إطار إلى
صخر مجهول الأصل.
\end{definition}

المخطّط الكلاسيكي هو مثلث Ti--Zr--Y لبيرس وكان (1973)، المرسوم للبازلتات. وهو
يستعمل عناصر قلّما يحرّكها التبدّل، فيمكن تطبيقه على حمم قديمة
متحوّلة ضاع إطارها. ويختبره الشكل أدناه على بازلتات المجموعة البالغ عددها 538
التي تتوفّر فيها العناصر الثلاثة. فمعظم بازلتات الجزر المحيطية يقع في
الحقل D داخل الصفائح، 152 من 173، ومعظم بازلت MORB في الحقل B، 142 من 164؛ لكن الحقل B
مشترك مع بازلتات الأقواس، و101 من أصل 201 بازلت أقواس تقع فيه أيضًا.
% ledger: geoplotr:tizry, data:b4-basalt-settings

\begin{omfigure}
\centering
\begin{tikzpicture}
\begin{axis}[width=0.62\linewidth, height=0.55\linewidth, axis equal image,
  hide axis, xmin=-0.04, xmax=1.04, ymin=-0.06, ymax=0.92, clip=false,
  legend style={font=\scriptsize, at={(1.0,0.95)}, anchor=north west}]
\addplot[only marks, mark=*, mark size=0.7pt, blue!70!black, opacity=0.7] table[x=x, y=y] {figdata/out/bachelor-2-series-tizry-morb.dat};
\addlegendentry{MORB}
\addplot[only marks, mark=square*, mark size=0.7pt, orange!85!black, opacity=0.7] table[x=x, y=y] {figdata/out/bachelor-2-series-tizry-oib.dat};
\addlegendentry{بازلت الجزر المحيطية}
\addplot[only marks, mark=triangle*, mark size=0.9pt, red!75!black, opacity=0.7] table[x=x, y=y] {figdata/out/bachelor-2-series-tizry-iab.dat};
\addlegendentry{بازلت الأقواس الجزرية}
\foreach \f in {A,B,C,D}
  \addplot[thick, black, forget plot] table[x=x, y=y] {figdata/out/bachelor-2-series-tizry-\f.dat};
\draw[thick] (axis cs:0,0) -- (axis cs:1,0) -- (axis cs:0.5,0.8660254) -- cycle;
\node[anchor=south] at (axis cs:0.5,0.866) {Ti/100};
\node[anchor=north east] at (axis cs:0,0) {Zr};
\node[anchor=north west] at (axis cs:1,0) {Y $\times$ 3};
\draw[thin] (axis cs:0.33,0.36) -- (axis cs:0.22,0.46) node[anchor=south east, font=\small\bfseries, inner sep=1pt] {D};
\draw[thin] (axis cs:0.36,0.12) -- (axis cs:0.24,0.06) node[anchor=north east, font=\small\bfseries, inner sep=1pt] {C};
\draw[thin] (axis cs:0.56,0.37) -- (axis cs:0.76,0.50) node[anchor=west, font=\small\bfseries, inner sep=1pt] {A};
\draw[thin] (axis cs:0.53,0.21) -- (axis cs:0.72,0.13) node[anchor=north west, font=\small\bfseries, inner sep=1pt] {B};
\end{axis}
\end{tikzpicture}
\par\smallskip{\small مخطّط Ti--Zr--Y لبيرس وكان (1973) مع البازلتات
المعروفة الإطار. A ثولييتات الأقواس الجزرية؛ B بازلت MORB وثولييتات الأقواس الجزرية
والبازلتات الكلسية القلوية؛ C البازلتات الكلسية القلوية؛ D بازلتات داخل الصفائح. Ti محسوب من
\ce{TiO2}؛ والعناصر الثلاثة بوحدة ppm.}
\end{omfigure}
% ledger: geoplotr:tizry, data:b4-basalt-settings

\begin{method}[استعمال مخطّط تمييز]\label{met:b2:magmatic-series:discrimination}
(1) تحقّق من أن الصخر من النوع الذي بُني له المخطّط، وهو عادةً بازلت،
وليس متراكمًا أو حممًا متطوّرة، تعكس عناصرها النزرة التبلورَ بقدر
ما تعكس المصدر. (2) فضّل العناصر غير المتحرّكة في التبدّل (Ti وZr وY وNb \omterm{def:b2:elements-in-the-earth:ree}{والعناصر
الأرضية النادرة} الثقيلة) للصخور المتحوّلة أو المتبدّلة. (3) ارسم مجموعة من العيّنات، لا عيّنة واحدة:
فنقطة مفردة قرب حدّ لا تحسم شيئًا. (4) اقرأ النتيجة احتمالًا:
فالحقول رُسمت حول معظم عيّنات مرجعها، لا كلها، وبعض الحقول
مشتركة. (5) أكّد بشواهد مستقلّة: مخطّطات أخرى، والنظائر،
والسياق الجيولوجي.
\end{method}

يمكن جعل الخطوة (4) كمّية (\cref{exo:b2:magmatic-series:8}). فإذا وقع بازلت في الحقل B، وكانت
الأطر الثلاثة متساوية الشيوع، فإن احتمال أن يكون بازلت MORB لا يتجاوز نحو 0.62.
والمصنِّفات الإحصائية المدرّبة على مثل هذه المجموعات ترث الحدّ نفسه: فالبازلت
يحمل ذاكرة مصدره وانصهاره، ويمكن لإطارين أن يصنعا صهارات
متشابهة.
% ledger: data:b4-basalt-settings

\begin{inthelab}[التحقّق من تحليل للصخر الكلي]
قبل رسم صخر على أيّ من هذه المخطّطات، يتحقّق علماء الصخور من التحليل. فالأكاسيد
الرئيسية، المقيسة مثلًا بفلورة الأشعة السينية، يجب أن يقارب مجموعها 100
مع المتطايرات: الماء وثاني أكسيد الكربون، المقيسين مباشرة أو معًا على هيئة
الفقد بالحرق، أي الكتلة التي يفقدها مسحوق عند التسخين الشديد. والفقد الكبير
بالحرق علامة صخر متبدّل، قد تكون قلوياته تحرّكت؛ ونظام مخطّط
TAS يرفض التحاليل التي فيها أكثر من \qty{2}{\percent} من الماء أو
\qty{0.5}{\percent} من ثاني أكسيد الكربون. ثم يُعاد حساب التحاليل إلى 100 خالية من
المتطايرات، ويُعبَّر عن الحديد بقيمة \ce{FeO} الكلية أو يُقسم بين الحديدوز والحديديك،
وتُقارن المواد المرجعية المحلَّلة في الدفعة نفسها بقيمها
المعتمدة (\cref{ch:b2:field-lab-methods-2}).
\end{inthelab}
% ledger: bgs:tas-use

\section{التمارين}

\begin{exercise}[$\star$]\label{exo:b2:magmatic-series:1}
سمِّ الصخور البركانية التالية على مخطّط TAS (\ce{SiO2} و\ce{Na2O} و\ce{K2O}
بوحدة \unit{\percent}، بعد استبعاد المتطايرات): (أ) 49.60، 2.22، 0.51 (بازلت هاواي
المرجعي BHVO-2)؛ (ب) 58.5، 3.9، 1.6؛ (ج) 72.0، 4.0، 4.2؛ (د) 47.0، 4.5، 2.5.
وفي (د)، سمِّ الصخرين الممكنين وقل ما الذي يحسم بينهما.
\end{exercise}
% ledger: usgs:bhvo2-major, bgs:tas

\begin{exercise}[$\star$]\label{exo:b2:magmatic-series:2}
تعطي حممٌ \ce{SiO2} 51.60، \ce{TiO2} 1.20، \ce{Al2O3} 16.10، \ce{FeO_T} 9.80،
\ce{MgO} 6.60، \ce{CaO} 8.20، \ce{Na2O} 3.50، \ce{K2O} 1.40، \ce{H2O} 1.40
و\ce{CO2} \qty{0.20}{\percent}. (أ) هل التحليل مقبول لمخطّط TAS؟
(ب) أعد حسابه خاليًا من المتطايرات. (ج) سمِّ الصخر قبل إعادة
الحساب وبعدها. (د) هل هو منخفض البوتاسيوم أم متوسطه أم عاليه؟
\end{exercise}
% ledger: bgs:tas-use, bgs:k2o-sio2

\begin{exercise}[$\star$]\label{exo:b2:magmatic-series:3}
سمِّ على \omterm{def:b2:magmatic-series:qapf}{مخطّط QAPF}: (أ) صخرًا خشن الحبيبات فيه \qty{25}{\percent} من الكوارتز،
و\qty{30}{\percent} من فلسبار القلويات، و\qty{35}{\percent} من البلاجيوكليز، و\qty{8}{\percent} من
البيوتيت و\qty{2}{\percent} من الهورنبلند؛ (ب) صخرًا فيه \qty{5}{\percent} من الكوارتز،
و\qty{5}{\percent} من فلسبار القلويات، و\qty{55}{\percent} من البلاجيوكليز
و\qty{35}{\percent} من الكلينوبيروكسين.
\end{exercise}
% ledger: bgs:qapf

\begin{exercise}[$\star$]\label{exo:b2:magmatic-series:4}
صنّف إلى منخفض البوتاسيوم أو متوسطه أو عاليه، بخطوط شكل \ce{K2O}--\ce{SiO2}:
(أ) 50.0، 0.40؛ (ب) 56.0، 1.50؛ (ج) 62.0، 2.90؛ (د) 66.0، 1.00 (\ce{SiO2} ثم
\ce{K2O}، بوحدة \unit{\percent}).
\end{exercise}
% ledger: bgs:k2o-sio2

\begin{exercise}[$\star\star$]\label{exo:b2:magmatic-series:5}
(أ) احسب $A$ و$F$ و$M$ لبازلت فيه \ce{FeO_T} 14.5، \ce{MgO} 4.5
و\ce{Na2O + K2O} \qty{3.2}{\percent}، ولأنديزيت فيه 6.1 و3.8 و5.05.
(ب) يمرّ الخط الفاصل بالنقاط $(F, A, M) = (54.5, 23.0, 22.5)$ و$(52.5, 30.0,
17.5)$ و$(49.6, 35.0, 15.4)$، وأولاها أغناها بالحديد؛
والجهة الثولييتية هي جهة F. في أيّ جهة يقع كل صخر؟ (ج) أيّ
معدن، إذ يتبلور مبكرًا، يُبقي الأنديزيت حيث هو؟
\end{exercise}
% ledger: geoplotr:afm, vp21:series

\begin{exercise}[$\star\star$]\label{exo:b2:magmatic-series:6}
لبازلت MORB الوسيط في المجموعة \qty{4}{ppm} من النيوبيوم و\qty{4.12}{ppm} من اللانثانوم،
ولبازلت الجزر المحيطية الوسيط 38 و30.8، ولبازلت الأقواس الوسيط 3 و6.1. ومع
قيمتي \omterm{def:b2:elements-in-the-earth:bse}{الوشاح البدائي}، 0.658 للنيوبيوم و\qty{0.648}{ppm} للانثانوم، احسب القيم المعايَرة
و$\mathrm{Nb_N}/\mathrm{La_N}$ لكلٍّ منها. أيّها فيه شذوذ نيوبيوم، ولماذا؟
\end{exercise}
% ledger: data:b4-basalt-settings, ms95:pm-trace

\begin{exercise}[$\star\star$]\label{exo:b2:magmatic-series:7}
يساوي Ti نسبة $0.5993$ من كتلة \ce{TiO2}. احسب إحداثيات الرسم على مخطّط Ti--Zr--Y
(Ti/100 وZr و$3\mathrm Y$، بعد إعادة التطبيع إلى 100) لثلاثة بازلتات: (أ) \ce{TiO2}
\qty{1.6}{\percent}، وZr 110، وY \qty{32}{ppm}؛ (ب) 2.8، 220، 29؛ (ج) 0.8، 50، 20. حدّد
مواضعها على الشكل وقل ما يسمح لك كل موضع باستنتاجه.
\end{exercise}
% ledger: geoplotr:tizry

\begin{exercise}[$\star\star$]\label{exo:b2:magmatic-series:8}
في اختبار Ti--Zr--Y في هذا الفصل، يضمّ الحقل B عدد 142 من 164 بازلت MORB، و4 من 173
بازلت جزر محيطية و101 من 201 بازلت أقواس. (أ) ما كسر كل مجموعة الواقع
في B؟ (ب) إذا كانت الأطر الثلاثة متساوية الشيوع، فما احتمال أن يكون
بازلت يقع في B بازلت MORB؟ (ج) ماذا يسمح لك الحقل B وحده أن تقول؟
\end{exercise}
% ledger: data:b4-basalt-settings

\begin{exercise}[$\star\star$]\label{exo:b2:magmatic-series:9}
ريوليت فيه \ce{Na2O} 5.5، \ce{K2O} 4.5 و\ce{Al2O3} \qty{10.5}{\percent} (قيم $M$ هي
61.98 و94.20 و\qty{101.96}{g/mol}). احسب نسبته الجزيئية $(\ce{Na2O +
K2O})/\ce{Al2O3}$. (أ) هل هو فوق قلوي؟ (ب) يحوي فلسبار القلويات أكسيدًا قلويًا
واحدًا لكل \ce{Al2O3}: فما كسر القلويات الذي لا يمكنه دخوله؟
\end{exercise}
% ledger: bgs:tas-use

\begin{exercise}[$\star\star\star$]\label{exo:b2:magmatic-series:10}
في نموذج قضية لوغاريتم النسب، خذ $\lambda_A = 0.4$ و$\lambda_M =
0.8$. (أ) احسب ميل الاتجاه من أجل $\lambda_F = 0.5$ ومن أجل $\lambda_F =
0.7$. (ب) لصخور آيسلندا ميل قدره $-1.10$ ولصخور الكاسكيد
$-0.72$: فما قيمة $\lambda_F$ التي يقتضيها كلٌّ منهما؟ (ج) أيّ المتتاليتين تفقد حديدها أسرع،
وأيّ معدن يفسّر ذلك؟ (د) ماذا يحدث للنسبة $F/M$ على امتداد كل اتجاه؟
\end{exercise}
% ledger: vp21:model

\begin{exercise}[$\star\star\star$]\label{exo:b2:magmatic-series:11}
معيار مبسّط. اعمل بمولات الأكاسيد. ضع \ce{K2O} كله في الأورثوكليز
(\ce{K2O.Al2O3.6SiO2})، و\ce{Na2O} كله في الألبيت (\ce{Na2O.Al2O3.6SiO2})،
و\ce{Al2O3} الباقي مع القدر نفسه من \ce{CaO} في الأنورثيت (\ce{CaO.Al2O3.2SiO2})،
و\ce{CaO} الباقي مع القدر نفسه من $(\ce{Mg,Fe})\ce{O}$ في الديوبسيد
(\ce{CaO.(Mg,Fe)O.2SiO2}) وبقية $(\ce{Mg,Fe})\ce{O}$ في الهيبرستين
(\ce{(Mg,Fe)O.SiO2}). فإن نقصت السيليكا، حوّل الهيبرستين إلى أوليفين
(\ce{2(Mg,Fe)O.SiO2})؛ وإن بقيت ناقصة، حوّل الألبيت إلى نيفلين
(\ce{Na2O.Al2O3.2SiO2}). طبّقه على (أ) BHVO-2: \ce{SiO2} 49.60، \ce{Al2O3} 13.44،
\ce{FeO_T} 11.15، \ce{MgO} 7.257، \ce{CaO} 11.40، \ce{Na2O} 2.219، \ce{K2O}
\qty{0.513}{\percent}؛ (ب) بازلت فيه 45.5 و15.0 و11.5 و8.5 و10.0 و3.8 و1.5 (الكتل
المولية: \ce{SiO2} 60.08، \ce{Al2O3} 101.96، \ce{FeO} 71.84، \ce{MgO} 40.30، \ce{CaO}
56.08، \ce{Na2O} 61.98، \ce{K2O} \qty{94.20}{g/mol}). هل كلٌّ منهما مشبّع بالسيليكا؟ وإلى
أيّ متتالية ينتمي كل بازلت؟
\end{exercise}
% ledger: usgs:bhvo2-major, bgs:series-terms

\begin{exercise}[$\star\star\star$]\label{exo:b2:magmatic-series:12}
في المجموعة، تتحقّق $\mathrm{Nb}/\mathrm{La} < 0.7$ في 116 من 141 بازلت أقواس،
و21 من 101 بازلت MORB و3 من 137 بازلت جزر محيطية حُلّلت للعنصرين. (أ) احسب
الكسور الثلاثة. (ب) مع نسب متساوية للأطر الثلاثة، ما
احتمال أن يكون بازلت فيه $\mathrm{Nb}/\mathrm{La} < 0.7$ بازلت أقواس؟ (ج)
أعد الحساب لإقليم فيه \qty{60}{\percent} من البازلتات بازلت MORB، و\qty{30}{\percent}
بازلت جزر محيطية و\qty{10}{\percent} بازلت أقواس. (د) ما الدرس الذي تحمله (ج)
لاستعمال مخطّطات التمييز؟
\end{exercise}
% ledger: data:b4-basalt-settings

\section{مسألة: بركان واحد وثلاث حمم}

\begin{problem}[{مسألة نهاية الأسبوع --- من ثلاثة تحاليل إلى إطار بركان
والبلّورات التي خلّفها}]
\label{pb:b2:magmatic-series:1}
تعطي ثلاث حمم من بركان طبقي (مختلَق)، من أقدم تدفّقات جوانبه إلى أحدث
قبابه، ما يلي (\unit{\percent}):

\smallskip\noindent{\small
\begin{tabular}{@{}lrrrrrrrrrr@{}}
 & \ce{SiO2} & \ce{TiO2} & \ce{Al2O3} & \ce{FeO_T} & \ce{MgO} & \ce{CaO} & \ce{Na2O} & \ce{K2O} & \ce{H2O} & \ce{CO2}\\
الحمم 1 & 50.4 & 0.95 & 17.6 & 8.6 & 7.9 & 10.2 & 2.9 & 0.62 & 0.70 & 0.10\\
الحمم 2 & 58.6 & 0.80 & 17.4 & 6.1 & 3.8 & 6.9 & 3.9 & 1.15 & 1.10 & 0.05\\
الحمم 3 & 65.3 & 0.55 & 16.2 & 4.0 & 1.6 & 4.3 & 4.4 & 1.85 & 1.20 & 0.05\\
\end{tabular}}

\smallskip\noindent وبوحدة ppm، للحمم 1 و2 و3 بالترتيب: الروبيديوم 10 و20 و32؛ الثوريوم 1.1 و2.1 و3.4؛ النيوبيوم 2.0 و3.6 و5.6؛
اللانثانوم 7 و13 و20؛ السترونشيوم 520 و480 و380؛ الزركونيوم 70 و125 و190؛ الإتريوم 18 و22 و24. \omterm{def:b2:elements-in-the-earth:bse}{والوشاح البدائي}: 0.658
للنيوبيوم، و0.648 للانثانوم، و\qty{0.0795}{ppm} للثوريوم.

\noindent\textbf{الجزء الأول --- تسمية الحمم.}
\begin{enumerate}
\item اجمع كل تحليل. هل التحاليل الثلاثة مقبولة لمخطّط TAS؟
\item أعد حساب الحمم 1 خاليةً من المتطايرات: أعطِ العامل، و\ce{SiO2}
و\ce{Na2O + K2O}.
\item افعل الشيء نفسه للحمم 2 و3.
\item سمِّ الحمم الثلاث على مخطّط TAS.
\item هل كانت الأسماء ستتغيّر دون إعادة الحساب؟
\item تبلور صخر خشن الحبيبات بتركيب الحمم 2 في جذور
البركان. أيّ نظام يسمّيه، وفي أيّ حقل تتوقّعه؟
\end{enumerate}

\noindent\textbf{الجزء الثاني --- أيّ متتالية؟}
\begin{enumerate}[resume]
\item صنّف كل حمم إلى منخفضة البوتاسيوم أو متوسطته أو عاليته.
\item احسب $A$ و$F$ و$M$ لكلٍّ منها (خاليةً من المتطايرات).
\item يمرّ الخط الفاصل بالنقاط $(F, A, M) = (42.8, 17.1, 40.1)$ و$(46.9, 18.1,
35.0)$ و$(52.5, 30.0, 17.5)$ و$(49.6, 35.0, 15.4)$ و$(40.8, 50.0, 9.2)$ و$(37.9,
55.0, 7.1)$، والحقل الثولييتي في جهة F. باستكمال $F$ خطيًا بدلالة
$A$ بين النقاط المتجاورة، في أيّ جهة تقع كل حمم؟
\item احسب \ce{FeO_T}/\ce{MgO} لكل حمم. هل هناك إثراء بالحديد؟
\item إلى أيّ متتالية ينتمي البركان؟ وأيّ معدن وأيّ ظروف
تفسّر ذلك؟
\end{enumerate}

\noindent\textbf{الجزء الثالث --- أيّ إطار؟}
\begin{enumerate}[resume]
\item عاير Nb وLa وTh في الحمم 1 إلى \omterm{def:b2:elements-in-the-earth:bse}{الوشاح البدائي}.
\item احسب $\mathrm{Nb_N}/\mathrm{La_N}$ للحمم 1 وقارنها بالقيم الوسيطة
في \cref{exo:b2:magmatic-series:6}.
\item احسب $\mathrm{Th_N}/\mathrm{Nb_N}$. ما الذي حمل الثوريوم؟
\item احسب إحداثيات Ti--Zr--Y للحمم 1 (بأكسيد \ce{TiO2} من التحليل الخالي من
المتطايرات). ماذا يقول حقلها؟ ولماذا لا ترسم الحمم 2 و3؟
\item استنتج: أين نما هذا البركان؟
\end{enumerate}

\noindent\textbf{الجزء الرابع --- كم تبلور؟}
\begin{enumerate}[resume]
\item لا يدخل البوتاس أيًّا من المعادن المبكرة. من \ce{K2O} الخالي من المتطايرات في
الحمم 1 و2، احسب كسر السائل $F$ الباقي حين تكوّنت الحمم 2، إن كانت
مشتقّة من الحمم 1 \omterm{def:b2:magmatic-differentiation:crystallisation}{بالتبلور التجزيئي}.
\item افعل الشيء نفسه بالروبيديوم وبالثوريوم.
\item للزركونيوم معامل كلي $D \approx 0.1$. استعمل قانون رايلي لتقدير $F$ منه.
\item ينخفض السترونشيوم من الحمم 2 إلى الحمم 3. احسب $F$ من البوتاس بين الحمم 2
و3، واستعمله لتقدير $D$ الكلي للسترونشيوم على ذلك المجال. أيّ معدن
مسؤول عن ذلك؟
\item من الحمم 1 إلى الحمم 3، ما كسر السائل الذي تبلور؟
\item لو استوعبت \omterm{def:b2:magma-emplacement:magma}{الصهارة} قشرة غرانيتية غنية بالبوتاس، هل سيكون تقدير
البوتاس مرتفعًا أكثر من اللازم أم منخفضًا أكثر من اللازم؟
\item خذ متوسط التقديرات الأربعة في الجزء الرابع للانتقال من الحمم 1 إلى 2 واذكر كتلة
البلّورات المتخلّفة لكل \qty{100}{kg} من البازلت لصنع الأنديزيت، وهي الرقم الأساسي
في هذه المسألة.
\end{enumerate}
\end{problem}
% ledger: bgs:tas-use, bgs:tas, bgs:k2o-sio2, geoplotr:afm, geoplotr:tizry, ms95:pm-trace
